% ==============================
% Acknowledgements
% ==============================

\chapter*{Acknowledgements}
\addcontentsline{toc}{chapter}{Acknowledgements}

I would like to express my sincere gratitude to all those who contributed to the successful completion of this Final Year Project.

First and foremost, I am deeply grateful to my \textbf{project supervisor} for their invaluable guidance, constructive feedback, and continuous support throughout this project. Their expertise in quantum computing and software engineering provided crucial insights that shaped the direction and quality of this work. Their patience in reviewing my code, suggesting improvements, and challenging me to think critically about design decisions has been instrumental in my growth as a developer and researcher.

I extend my heartfelt thanks to my \textbf{family} for their unwavering support, encouragement, and understanding during the demanding months of this project. Their belief in my abilities and their willingness to accommodate my long hours of work provided the emotional foundation necessary to persevere through challenges.

I am grateful to my \textbf{fellow students and colleagues} who participated in user acceptance testing, provided valuable feedback on the system's usability, and engaged in stimulating discussions about quantum computing concepts. Their diverse perspectives helped identify areas for improvement and validated the educational value of the Quantum Virtual Machine.

I would like to acknowledge the \textbf{open-source quantum computing community}, particularly the developers of Qiskit, ProjectQ, Cirq, and PyQuil, whose pioneering work in quantum software development provided inspiration and technical reference for this project. The availability of research papers, documentation, and open-source implementations accelerated my understanding of quantum compilation techniques and best practices.

Special thanks to the creators and maintainers of the \textbf{software libraries and tools} that made this project possible: NumPy for efficient numerical computation, Lark for robust parsing, FastAPI for modern web services, Matplotlib for visualization, pytest for comprehensive testing, and the Python community for creating an excellent ecosystem for scientific computing.

I am grateful to my \textbf{university and department} for providing the resources, infrastructure, and academic environment that enabled this research. Access to computational resources, library facilities, and technical support services was essential for the successful completion of this project.

Finally, I would like to thank the \textbf{pioneers of quantum computing}—from Richard Feynman and David Deutsch who envisioned quantum computation, to Peter Shor and Lov Grover who demonstrated its potential, to the countless researchers and engineers working today to make quantum computing a practical reality. Their groundbreaking work continues to inspire new generations of quantum computing enthusiasts.

This project represents not just my individual effort, but the collective support, knowledge, and encouragement of all these individuals and communities. Any errors or shortcomings in this work remain entirely my own responsibility.

\vspace{1cm}

\begin{flushright}
\textit{[Muhammad Abdul Qayyoum]}\\
\textit{April 2026}
\end{flushright}
